\clearpage

\newcommand{\evalsystemprompts}[1]{%
\textbf{Shared system instruction.} You are an evaluator of expert-domain
scientific writing. You will get a query--answer pair along with criteria
explaining the specific evaluation aspect and the scoring rubric. You should
evaluate whether the answer satisfy the query based on the given criteria.

\medskip
\hypertarget{#1-cot}{\textbf{CoT format.}} First output your reasoning enclosed
between \texttt{<reasoning>} and \texttt{</reasoning>}. Then output your score
enclosed between \texttt{<score>} and \texttt{</score>}. Inside
\texttt{<score>} provide only the numeric score and nothing else.

\medskip
\hypertarget{#1-label-only}{\textbf{Label-only format.}} Output your score
enclosed between \texttt{<score>} and \texttt{</score>}. Inside
\texttt{<score>} provide only the numeric score and nothing else. Do not output
any reasoning or other text.}

\begin{figure*}[p]
\centering
\fbox{\begin{minipage}{0.95\textwidth}
\raggedright\footnotesize
\textbf{Related Work--Coherence}

\medskip
\evalsystemprompts{prompt-rw-coherence}

\medskip
\textbf{User prompt.}

\texttt{[QUERY]:} Given a context from a scientific paper, your task is to
write a coherent citation sentence that cites the given paper with a specific
citation number. The citation number will also be provided with the context.

\medskip
\texttt{[CRITERIA]:} The paper context supports (entails) the citation
sentence. When a sentence references multiple papers, the given paper need
only support the part associated with its citation number. Scoring rubric:

0: The context does not align with the claim associated with the citation
number or with the overall meaning of the citation sentence.

1: The context is compatible with the claim associated with the citation
number.

\medskip
\texttt{[CONTEXT]:} This paper proposes a 3D shape descriptor network, which
is a deep convolutional energy-based model for modeling volumetric shape
patterns. Maximum-likelihood training follows an ``analysis by synthesis''
scheme and can be interpreted as mode seeking and mode shifting. The model can
synthesize 3D shapes through MCMC sampling, train a 3D generator through MCMC
teaching, and support 3D object recovery and super-resolution. Experiments
show that it generates realistic shapes and supports 3D shape analysis.

\medskip
\texttt{CITATION NUMBER:} 5

\medskip
\texttt{[ANSWER]:} Recent work has introduced the EBM with a ConvNet
potential [4], [5].
\end{minipage}}
\caption{Matched CoT and label-only evaluation prompts for Related
Work--Coherence. The task input is shared across the two formats.}
\label{fig:prompt-rw-coherence}
\end{figure*}

\clearpage

\begin{figure*}[p]
\centering
\fbox{\begin{minipage}{0.95\textwidth}
\raggedright\footnotesize
\textbf{Related Work--Positioning Type}

\medskip
\evalsystemprompts{prompt-rw-positioning-type}

\medskip
\textbf{User prompt.}

\texttt{[QUERY]:} Your task is to write a related-work section for a
scientific paper that states the paper's contribution or its position among
the literature in each paragraph.

\medskip
\texttt{[CRITERIA]:} In a multi-paragraph related-work section, every
paragraph should state the paper's contribution and/or position in the
literature. The contribution or positioning statement should align with the
specific focus of its paragraph. Scoring rubric:

0: The answer fails to state the paper's contribution and/or position in each
related-work paragraph.

1: The answer states the paper's contribution and/or position in each
related-work paragraph.

\medskip
\texttt{[ANSWER]:} In several machine-learning settings, negative samples are
produced from a statistical generative model. [1] generates negative data
using GANs for semi-supervised learning and novelty detection. [2] proposes a
theoretical framework that relies on an oracle classifying a sample as valid
or invalid, without providing a practical implementation. [3] uses
adversarial training to generate hard negatives for NLP tasks. [4] employs a
GAN to learn the negative-data distribution for positive-unlabeled learning.

In contrastive unsupervised learning, negative examples are assumed to be
semantically farther from positive samples. Word2Vec [5] treats samples from a
different context as negatives, while CPC-based methods [6], such as momentum
contrast [7], use augmentations of different images. Negative data
augmentation generates negatives by transforming existing samples, thereby
incorporating useful inductive biases [8].
\end{minipage}}
\caption{Matched CoT and label-only evaluation prompts for Related
Work--Positioning Type. The task input is shared across the two formats.}
\label{fig:prompt-rw-positioning-type}
\end{figure*}

\clearpage

\begin{figure*}[p]
\centering
\fbox{\begin{minipage}{0.95\textwidth}
\raggedright\footnotesize
\textbf{Related Work--Positioning Check}

\medskip
\evalsystemprompts{prompt-rw-positioning-check}

\medskip
\textbf{User prompt.}

\texttt{[QUERY]:} Your task is to write a related-work paragraph for a
scientific paper that states the paper's contribution or its position among
the literature.

\medskip
\texttt{[CRITERIA]:} The generated paragraph should explicitly or implicitly
mention the main paper's contribution or position among existing literature.
The contribution and/or positioning statement should align with the specific
focus of the paragraph. Scoring rubric:

0: The paragraph fails to explicitly or implicitly mention the main paper's
contribution or position among existing literature.

1: The paragraph explicitly or implicitly mentions the main paper's
contribution or position among existing literature.

\medskip
\texttt{[ANSWER]:} Another line of research focuses on transforming extracted
feature spaces into representations that are more discriminative for novel
classes. Approaches that learn such transformations include metric-learning
techniques and feature-space adaptation methods [9--12].
\end{minipage}}
\caption{Matched CoT and label-only evaluation prompts for Related
Work--Positioning Check. The task input is shared across the two formats.}
\label{fig:prompt-rw-positioning-check}
\end{figure*}

\clearpage

\begin{figure*}[p]
\centering
\fbox{\begin{minipage}{0.95\textwidth}
\raggedright\footnotesize
\textbf{RevUtil--Actionability}

\medskip
\evalsystemprompts{prompt-ru-actionability}

\medskip
\textbf{User prompt.}

\texttt{[QUERY]:} Your task is to write a review comment for a scientific
paper. The comment should be actionable. The actions should be clearly
identifiable and concrete.

\medskip
\texttt{[CRITERIA]:} Explicit actions are direct or apparent, whereas implicit
actions must be inferred. Concrete actions tell authors exactly what to change
and how to implement the change; vague actions do not. Clarification questions
count as explicit when they specify a direct action. Scoring rubric:

1: The comment lacks meaningful information that could help the authors
improve the paper.

2: The comment implies an action, but the action is vague and lacks
implementation details.

3: The comment explicitly states an action but is vague about how to execute
it.

4: The comment implies an action and concretely explains how to implement it.

5: The comment explicitly states an action and provides concrete details for
implementing it.

\medskip
\texttt{[ANSWER]:} 5. Why is TS-CAN better than MTTS-CAN in some cases?
\end{minipage}}
\caption{Matched CoT and label-only evaluation prompts for
RevUtil--Actionability. The task input is shared across the two formats.}
\label{fig:prompt-ru-actionability}
\end{figure*}

\clearpage

\begin{figure*}[p]
\centering
\fbox{\begin{minipage}{0.95\textwidth}
\raggedright\footnotesize
\textbf{RevUtil--Grounding Specificity}

\medskip
\evalsystemprompts{prompt-ru-grounding}

\medskip
\textbf{User prompt.}

\texttt{[QUERY]:} Your task is to write a review comment for a scientific
paper. The comment should refer to a specific part of the paper and clearly
identify the issue with that part.

\medskip
\texttt{[CRITERIA]:} A fully grounded comment allows the author to pinpoint
the section, table, figure, or unique aspect being addressed. A specific
comment explains what is wrong or missing in that part and provides concrete
examples when external work is mentioned. Scoring rubric:

1: The comment is neither grounded nor specific.

2: The addressed part cannot be identified confidently, and the issue is not
specified.

3: The addressed part cannot be identified confidently, but the issue is
specified clearly.

4: The addressed part is explicit or obvious, but the issue is not specified.

5: The addressed part is explicit or obvious, and the issue is specified.

\medskip
\texttt{[ANSWER]:} 4. Figure 3 lacks a description in the main body.
\end{minipage}}
\caption{Matched CoT and label-only evaluation prompts for RevUtil--Grounding
Specificity. The task input is shared across the two formats.}
\label{fig:prompt-ru-grounding}
\end{figure*}

\clearpage

\begin{figure*}[p]
\centering
\fbox{\begin{minipage}{0.95\textwidth}
\raggedright\footnotesize
\textbf{RevUtil--Helpfulness}

\medskip
\evalsystemprompts{prompt-ru-helpfulness}

\medskip
\textbf{User prompt.}

\texttt{[QUERY]:} Your task is to write a review comment for a scientific
paper. The comment should be useful for helping the authors improve the paper.

\medskip
\texttt{[CRITERIA]:} A helpful review should be actionable, grounded in a
specific part of the paper, and provide justification or evidence for its
claims. Scoring rubric:

1: The comment identifies no meaningful weakness or improvement and provides
no actionable feedback.

2: The comment identifies an issue but is vague, unclear, or minimally useful.

3: The comment identifies an issue or improvement but is incomplete or lacks
depth.

4: The comment provides clear and actionable feedback but could be more
comprehensive.

5: The comment thoroughly identifies weaknesses and provides detailed,
actionable, and constructive suggestions.

\medskip
\texttt{[ANSWER]:} 1) Which model is better: HMM/DNN, CTC, or Enc-Dec?
\end{minipage}}
\caption{Matched CoT and label-only evaluation prompts for
RevUtil--Helpfulness. The task input is shared across the two formats.}
\label{fig:prompt-ru-helpfulness}
\end{figure*}

\clearpage

\begin{figure*}[p]
\centering
\fbox{\begin{minipage}{0.95\textwidth}
\raggedright\footnotesize
\textbf{RevUtil--Verifiability}

\medskip
\evalsystemprompts{prompt-ru-verifiability}

\medskip
\textbf{User prompt.}

\texttt{[QUERY]:} Your task is to write a review comment for a scientific
paper. The claims in the comment should be justified through logical
reasoning, common sense, or external sources.

\medskip
\texttt{[CRITERIA]:} A claim may be justified through supporting logic,
common knowledge in the field, or relevant external references. Scoring
rubric:

1: The comment contains a claim without supporting evidence or justification.

2: The comment provides vague, insufficient, or incomplete support.

3: The comment provides support but omits important examples, explanations,
or references.

4: The claim is sufficiently supported but has minor gaps.

5: The claim is thoroughly supported by explicit and robust reasoning,
references, data, or common-sense arguments.

\medskip
\texttt{[ANSWER]:} 4. It is probably better to add a conclusion section.
\end{minipage}}
\caption{Matched CoT and label-only evaluation prompts for
RevUtil--Verifiability. The task input is shared across the two formats.}
\label{fig:prompt-ru-verifiability}
\end{figure*}
